\newcommand{\MercurySymbol}{\text{\mdseries\AstroSymbols ☿}}
\bigskip\Needspace{10\baselineskip}
\pdfbookmark[1]{Transit of Mercury Across the Sun, 3 November 1802}{mercur18021103}
\begin{center}
TRANSIT OF MERCURY ACROSS THE SUN, 3 NOVEMBER 1802.\par
\medskip\hrule\medskip
Monatliche Correspondenz. Vol. VI. pp.\ 567—575. December 1802.\par
\medskip\hrule
\end{center}\medskip
—\quad —\quad —

In \textit{Braunschweig}, Dr.\ \textsc{Gauss} observed the \textit{internal} contact of egress at $0^{\mathrm u}\ 39^{\mathrm m}\ 16^{\mathrm s}$ and \textit{external} contact at $0^{\mathrm u}\ 40^{\mathrm m}\ 48^{\mathrm s}$ true time. [Manuscript. Egress of $\MercurySymbol$ observed in Braunschweig on 9 Nov.\ $0^{\mathrm u}\ 23^{\mathrm m}\ 17^{\mathrm s}$ internal contact, $0^{\mathrm u}\ 24^{\mathrm m}\ 49^{\mathrm s}$ external contact, mean time.] He used a two-foot achromat made by \textsc{Baumann}, in which Mercury and the beautiful group of sunspots it passed by could be seen very clearly. Unfortunately, he was unable to determine the time as accurately as he would have liked. Dr.\ \textsc{Gauss} shares this difficulty
\SourceLeaf{237}%
at this time of year with several astronomers who have no transit instrument. As it was, he had to make do with a few individual altitudes of the Sun, taken after 1 o’clock on an uncovered oil horizon, to determine the time. For this reason, and because this is the first observation of its kind he has made, he asks for indulgence if his observation should fail to agree exactly with others. Besides observing the egress, he also tried using the telescope as a ring micrometer and measured several differences in right ascension between Mercury and the center of the Sun. He himself has little confidence in their accuracy. They can in any case be checked by calculation.
\begin{center}
\begin{tabular}{r|r}
True time&\shortstack{Difference in RA\\$\odot-\MercurySymbol$}\\\hline
$23^{\mathrm u}\ 26^{\mathrm m}\ 6^{\mathrm s}$&$7'\ 51''$\\
$23\ 42\ 30$&$9\ 29$\\
$23\ 50\ 28$&$10\ 3$\\
$0\ 3\ 25$&$11\ 37$\\
$0\ 25\ 14$&$13\ 41$
\end{tabular}
\end{center}
—\quad —\quad —
\begin{flushright}[\textsc{von Zach.}]\end{flushright}

